\section{Introduction}
\label{sec:intro}

A technical report earns its keep by being checkable: the reader follows a number to an
experiment, and the experiment to a line of code. Typography does not create that value,
but it decides whether anyone keeps reading, and whether the author re-derives the same
layout decisions for every new report.

Three things kept costing us time. First, every blank document meant re-tuning margins,
font sizes, caption punctuation and table rules, and no two reports came out alike.
Second, mixed Chinese and English kept tripping over engines: Chinese needs XeLaTeX for
line breaking, English is cheapest with pdfLaTeX, and both must look like the same
document. Third, numbers were written without a convention, so a table could contain
both \texttt{62.4} and \texttt{62.40} and the reader had no way to tell which was a typo.

This template addresses those three problems. Our contributions are:

\begin{contributions}
  \item a single class file, \texttt{boogonreport.cls}, that carries every layout
        decision: title block, abstract, captions, three-rule tables, code listings,
        references and appendices;
  \item one \texttt{zh} option that switches both the font pipeline and the
        localised strings, replacing a page of \verb|\renewcommand| boilerplate;
  \item build scripts and CI that produce the same PDF from any machine with TeX~Live.
\end{contributions}

The report is organised as follows: \secref{sec:layout} covers the page geometry,
\secref{sec:type} the two font pipelines, \secref{sec:figure} the rules for figures,
tables and numbers, \secref{sec:build} the build and release flow, and
\secref{sec:limit} what the template does not attempt.
